% !TeX program = xelatex
% 自编教学资料；A3 张贴用，含例题、练习及解答，不作为无答案试卷。
\documentclass[UTF8,fontset=windows,zihao=-4]{ctexart}
\usepackage[a3paper,margin=17mm,headheight=18pt,headsep=8mm,footskip=9mm]{geometry}
\usepackage{amsmath,amssymb,multicol,array,booktabs,tikz,fancyhdr}
\setlength{\parindent}{0pt}
\setlength{\parskip}{3pt}
\setlength{\columnsep}{12mm}
\setlength{\columnseprule}{0.3pt}
\pagestyle{fancy}\fancyhf{}
\fancyhead[L]{\small 方程与不等式基础知识}
\fancyhead[R]{\small 实数范围}
\fancyfoot[C]{\small\thepage}
\newcommand{\sheet}[1]{\begin{center}\heiti\fontsize{24}{30}\selectfont #1\end{center}\vspace{3mm}}
\newcommand{\topic}[1]{\par\vspace{7pt}{\heiti\fontsize{17}{23}\selectfont #1}\par\vspace{2pt}}
\newcommand{\ex}[1]{\par\textbf{例}\quad #1\par}
\newcommand{\cross}[4]{%
\begin{tikzpicture}[baseline=(current bounding box.center),x=1cm,y=1cm]
\node (a) at (0,0.8) {$#1$};\node (b) at (1.7,0.8) {$#2$};
\node (c) at (0,-0.1) {$#3$};\node (d) at (1.7,-0.1) {$#4$};
\draw (a.south east)--(d.north west);\draw (c.north east)--(b.south west);
\end{tikzpicture}}
\begin{document}
\raggedcolumns
\linespread{1}\fontsize{14}{19}\selectfont
\setlength{\abovedisplayskip}{5pt}\setlength{\belowdisplayskip}{5pt}
\setlength{\abovedisplayshortskip}{3pt}\setlength{\belowdisplayshortskip}{3pt}
\sheet{一、基本运算与一元二次方程}
\begin{multicols*}{2}
\topic{1. 字母、代数式与方程的解}
$2x$ 表示 $2\times x$，$x^2$ 表示 $x\times x$。例如 $x=-3$ 时，
\[2x=-6,\qquad x^2=(-3)^2=9.\]
使方程左右两边相等的未知数的值，叫作方程的\textbf{解}，也叫作\textbf{根}。

例如 $x+2=5$ 的解是 $x=3$，因为 $3+2=5$。检验一个数是不是解，可以把它代回原方程。
\topic{2. 负数运算与去括号}
同号相乘得正，异号相乘得负。
\[(-3)(-2)=6,\qquad(-3)\times2=-6.\]
注意平方与负号的运算顺序：
\[(-3)^2=9,\qquad -3^2=-(3^2)=-9.\]
去括号时，括号前的数要乘到每一项：
\[2(x-3)=2x-6,\quad -2(x-3)=-2x+6.\]
相同字母、相同次数的项才能合并：
\[3x+2x=5x.\]
$x^2$ 与 $2x$ 不是同类项，不能合并成一项。
\topic{3. 等式性质与移项}
等式两边同加、同减同一个数，或同乘、同除以同一个非零数，所得等式仍成立。

移项要改变该项的符号。例如
\[x+3=7\quad\Rightarrow\quad x=7-3=4.\]
这相当于原方程两边同时减去 $3$。
\ex{解方程 $2(x-1)+3=7$。}
\[
\begin{aligned}
2x-2+3&=7&&\text{（去括号）},\\
2x+1&=7&&\text{（合并）},\\
2x&=6&&\text{（两边减去 1）},\\
x&=3&&\text{（两边除以 2）}.
\end{aligned}
\]
\topic{4. 一元二次方程的一般形式}
\[\boxed{ax^2+bx+c=0\quad(a\ne0)}\]
$a,b,c$ 分别是二次项系数、一次项系数和常数项，\textbf{符号也包括在系数中}。

例如 $2x^2=3x+2$，整理为
\[2x^2-3x-2=0,\]
所以 $a=2,b=-3,c=-2$。方程 $x^2-5=0$ 中，$a=1,b=0,c=-5$。
\columnbreak
\topic{5. 直接开平方法}
实数的平方不小于 $0$。因此
\[
x^2=k\ \Longrightarrow\ \begin{cases}
x=\pm\sqrt{k},&k>0,\\
x=0,&k=0,\\
\text{无实数解},&k<0.
\end{cases}
\]
$\sqrt9=3$，但 $x^2=9$ 的根为 $3$ 和 $-3$。符号 $\pm$ 表示正、负两种情况。
\ex{解方程 $(x-1)^2=9$。}
把 $x-1$ 看作一个整体，开平方得
\[x-1=3\quad\text{或}\quad x-1=-3.\]
分别解得 $x=4$ 或 $x=-2$。
\ex{解方程 $(x+2)^2=5$。}
\[x+2=\pm\sqrt5,\qquad x=-2\pm\sqrt5.\]
$\sqrt5$ 是平方等于 $5$ 的正数，约为 $2.236$；答案一般保留根号。
\topic{6. 因式分解法}
若 $AB=0$，则 $A=0$ 或 $B=0$。解题时先将方程一边化为 $0$，再分解另一边。
\ex{解方程 $x^2=3x$。}
\[x^2-3x=0\quad\Rightarrow\quad x(x-3)=0.\]
所以 $x=0$ 或 $x-3=0$，即 $x=0$ 或 $x=3$。

\textbf{注意：}不能直接两边除以 $x$，因为 $x$ 可能为 $0$，这样会漏根。

常用的因式分解公式：
\[
\begin{aligned}
x^2-d^2&=(x-d)(x+d),\\
x^2+2dx+d^2&=(x+d)^2,\\
x^2-2dx+d^2&=(x-d)^2.
\end{aligned}
\]
例如 $x^2-9=(x-3)(x+3)$，所以方程 $x^2-9=0$ 的根为 $3$ 和 $-3$。
\topic{练习}
（1）解方程 $3x-4=8$。

（2）解方程 $(x+2)^2=16$。

（3）解方程 $x^2-4x=0$。
\topic{解答}
（1）$3x=12$，所以 $x=4$。

（2）$x+2=\pm4$，所以 $x=2$ 或 $-6$。

（3）$x(x-4)=0$，所以 $x=0$ 或 $4$。
\end{multicols*}
\newpage
\sheet{二、十字相乘法}
\begin{multicols*}{2}
\topic{1. 十字相乘是在检查什么？}
十字相乘法是分解二次三项式的一种方法。它的依据是乘法展开式：
\[(px+q)(rx+s)=prx^2+(ps+qr)x+qs.\]
要分解 $ax^2+bx+c$，就要找到 $p,q,r,s$，使
\[pr=a,\qquad qs=c,\qquad ps+qr=b.\]
\begin{center}\cross{px}{q}{rx}{s}\end{center}
左边竖着相乘得到 $ax^2$，右边竖着相乘得到 $c$；\textbf{交叉相乘后相加，必须得到一次项 $bx$}。检查通过后，按两行分别写成两个因式。
\topic{2. 二次项系数为 1 的例子}
\ex{分解 $x^2-5x+6$，并解方程 $x^2-5x+6=0$。}
先把 $x^2$ 拆为 $x\cdot x$；把 $6$ 拆为两个数的乘积。因为一次项是 $-5x$，试用 $-2$ 与 $-3$：
\begin{center}\cross{x}{-2}{x}{-3}\end{center}
交叉相乘，再相加：
\[x\cdot(-3)+x\cdot(-2)=-3x-2x=-5x.\]
与原式的一次项相同，所以
\[x^2-5x+6=(x-2)(x-3).\]
原方程变为 $(x-2)(x-3)=0$，于是
\[x-2=0\ \text{或}\ x-3=0,\]
所以 $x=2$ 或 $x=3$。
\topic{3. 怎样选择右边两个数的符号？}
对 $x^2+bx+c$，要找的两个数的\textbf{积为 $c$、和为 $b$}。

$c>0$：两个数同号。$b>0$ 时都为正，$b<0$ 时都为负。

$c<0$：两个数异号。它们的和的符号，由绝对值较大的那个数决定。

例如分解 $x^2+x-6$，可试 $3$ 与 $-2$，因为
\[3(-2)=-6,\qquad3+(-2)=1,\]
所以 $x^2+x-6=(x+3)(x-2)$。
\columnbreak
\topic{4. 二次项系数不是 1 的例子}
\ex{分解 $2x^2-3x-2$，并解方程 $2x^2-3x-2=0$。}
把 $2x^2$ 拆为 $2x\cdot x$；把 $-2$ 拆为 $1\cdot(-2)$。安排如下：
\begin{center}\cross{2x}{1}{x}{-2}\end{center}
交叉相乘，再相加：
\[2x\cdot(-2)+x\cdot1=-4x+x=-3x.\]
因此
\[2x^2-3x-2=(2x+1)(x-2).\]
原方程变为 $(2x+1)(x-2)=0$，所以
\[2x+1=0\quad\text{或}\quad x-2=0,\]
即 $x=-\dfrac12$ 或 $x=2$。

\textbf{为什么右边两个数的位置不能随意放？}

如果写成下面的排列：
\begin{center}\cross{2x}{-2}{x}{1}\end{center}
交叉乘积之和是 $2x-2x=0$，不是 $-3x$，所以这个排列不符合原式。
\topic{5. 分解后的检查}
把两个因式展开，看三项是否都与原式一致：
\[
\begin{aligned}
(2x+1)(x-2)
&=2x^2-4x+x-2\\
&=2x^2-3x-2.
\end{aligned}
\]
\textbf{注意：}十字相乘不是每次都能用小整数试出来。例如 $x^2+x-1$ 不能分解为两个整系数的一次因式，可改用配方法或公式法。
\topic{练习}
用十字相乘法分解左边，再解方程。

（1）$x^2-7x+12=0$。

（2）$x^2+x-6=0$。

（3）$2x^2+x-3=0$。
\topic{解答}
（1）$(x-3)(x-4)=0$，$x=3$ 或 $4$。

（2）$(x+3)(x-2)=0$，$x=-3$ 或 $2$。

（3）$(2x+3)(x-1)=0$，$x=-\dfrac32$ 或 $1$。交叉检查：$-2x+3x=x$。
\end{multicols*}
\newpage
\sheet{三、配方法、公式法与根的关系}
\begin{multicols*}{2}
\topic{1. 配方法}
把方程整理成“一个式子的平方等于一个数”，再开平方求解。

二次项系数先化为 $1$；移去常数项后，两边同加\textbf{一次项系数的一半的平方}。

例如 $x^2+4x$ 要补 $\left(\dfrac42\right)^2=4$，因为
\[x^2+4x+4=(x+2)^2.\]
\ex{解方程 $x^2+4x-1=0$。}
\[
\begin{aligned}
x^2+4x&=1,\\
x^2+4x+4&=1+4,\\
(x+2)^2&=5,\\
x+2&=\pm\sqrt5.
\end{aligned}
\]
所以 $x=-2+\sqrt5$ 或 $x=-2-\sqrt5$。

\textbf{注意：}只能两边同加，不能只改变左边。
\topic{2. 判别式与求根公式}
对 $ax^2+bx+c=0\ (a\ne0)$，令
\[\Delta=b^2-4ac.\]
“$\Delta$”读作德尔塔，叫作判别式。

$\Delta>0$：有两个不相等的实数根。

$\Delta=0$：有两个相等的实数根，即重根，只有一个不同的根值。

$\Delta<0$：没有实数根。

当 $\Delta\ge0$ 时，求根公式为
\[\boxed{x=\frac{-b\pm\sqrt{b^2-4ac}}{2a}}.\]
整个分子都要除以 $2a$，不是只除根号部分。
\ex{解方程 $2x^2-3x-2=0$。}
$a=2,b=-3,c=-2$，所以
\[\Delta=(-3)^2-4\times2\times(-2)=25.\]
代入求根公式：
\[x=\frac{3\pm5}{4},\]
所以 $x=2$ 或 $x=-\dfrac12$。
\columnbreak
\topic{3. 怎样选择解法？}
能整理成 $(x-h)^2=k$ 的，直接开平方。

各项有公因式的，先提公因式；容易分解的，可以用平方差、完全平方公式或十字相乘法。

一时看不出怎样分解的，可以用配方法或公式法。\textbf{十字相乘法属于因式分解法}，不是求根公式的必要步骤。
\topic{4. 韦达定理}
若 $ax^2+bx+c=0\ (a\ne0)$ 有两个实数根 $x_1,x_2$（可以相等），则
\[\boxed{x_1+x_2=-\frac ba,\qquad x_1x_2=\frac ca}.\]
即两根之和为 $-b/a$，两根之积为 $c/a$。使用实数根时，应先确认 $\Delta\ge0$。

这是因为
\[
\begin{aligned}
a(x-x_1)(x-x_2)
&=ax^2-a(x_1+x_2)x\\
&\quad+ax_1x_2,
\end{aligned}
\]
与原二次三项式比较系数，就得到上述关系。
\ex{不求根，求 $x^2-5x+3=0$ 两根的平方和。}
$\Delta=25-12=13>0$，两实根存在。
\[x_1+x_2=5,\qquad x_1x_2=3.\]
利用完全平方公式变形：
\[
\begin{aligned}
x_1^2+x_2^2&=(x_1+x_2)^2-2x_1x_2\\
&=25-6=19.
\end{aligned}
\]
\topic{练习}
（1）解方程 $x^2+6x+5=0$。

（2）用公式法解 $x^2-2x-3=0$。

（3）已知 $x^2-5x+m=0$ 的一个根为 $2$，求另一个根及 $m$。
\topic{解答}
（1）配方得 $(x+3)^2=4$，所以 $x=-1$ 或 $-5$。

（2）$\Delta=16$，$x=\dfrac{2\pm4}{2}$，所以 $x=3$ 或 $-1$。

（3）代入已知根得 $4-10+m=0$，故 $m=6$。方程化为 $(x-2)(x-3)=0$，另一个根是 $3$；也可用根的和算得 $5-2=3$。
\end{multicols*}
\newpage
\sheet{四、一元一次不等式与一元二次不等式}
\begin{multicols*}{2}
\topic{1. 不等式的解与基本性质}
解不等式，是求使不等式成立的所有未知数的值。例如 $x+2<5$ 的解为 $x<3$，包括满足条件的整数、小数等所有实数。

两边同加、同减同一个数，不等号方向不变；同乘、同除以同一个正数，方向不变；\textbf{同乘、同除以同一个负数，方向改变}。
\ex{解不等式 $3-2x\ge7$。}
\[-2x\ge4\quad\Rightarrow\quad x\le-2.\]
最后一步除以 $-2$，所以不等号改变方向。
\topic{2. 解集的表示}
\[
\begin{array}{c|c}
\text{不等式表示}&\text{区间表示}\\\hline
2<x<5&(2,5)\\
2\le x\le5&[2,5]\\
x\le-2&(-\infty,-2]\\
x<2\ \text{或}\ x>5&(-\infty,2)\cup(5,+\infty)
\end{array}
\]
圆括号不含端点，方括号包含端点。$\infty$ 表示无穷，无穷处只用圆括号；$\cup$ 表示把两部分合在一起。所有实数记为 $\mathbb R$，无解时解集为空集 $\varnothing$。
\topic{3. 二次三项式的正负}
以 $x^2-7x+10=(x-2)(x-5)$ 为例，两个根 $2,5$ 将数轴分为三个区间。
\[\renewcommand{\arraystretch}{1.15}
\begin{array}{c|ccc}
x&x<2&2<x<5&x>5\\\hline
x-2&-&+&+\\
x-5&-&-&+\\\hline
(x-2)(x-5)&+&-&+
\end{array}
\]
例如 $x<2$ 时，两个因式都为负，乘积为正。因此
\[
\begin{aligned}
x^2-7x+10<0&\Rightarrow 2<x<5,\\
x^2-7x+10>0&\Rightarrow x<2\ \text{或}\ x>5.
\end{aligned}
\]
若题目为 $\le0$，则包含两个端点，得 $[2,5]$。

\textbf{适用条件：}二次项系数为正，且有两个不同实根时，两根之间为负、两侧为正。端点处等于 $0$。
\columnbreak
\topic{4. 解题举例}
\ex{解不等式 $x^2-x\le6$。}
移项得 $x^2-x-6\le0$。对应方程为
\[(x-3)(x+2)=0,\]
两根是 $-2,3$。二次项系数为正，取两根之间并包含端点，解集为 $[-2,3]$。
\ex{解不等式 $-x^2+x+6>0$。}
两边乘 $-1$，注意改变不等号方向：
\[x^2-x-6<0.\]
两根仍为 $-2,3$，取两根之间但不含端点，解集为 $(-2,3)$。
\topic{5. 重根与无实根的情况}
以下均已将二次项系数化为正。

\textbf{（1）$\Delta=0$。}例如
\[x^2-4x+4=(x-2)^2\ge0.\]
只有 $x=2$ 时取到 $0$，所以
\[
\begin{array}{c|c}
\text{不等式}&\text{解}\\\hline
(x-2)^2>0&x\ne2\\
(x-2)^2\ge0&\text{所有实数}\\
(x-2)^2<0&\text{无解}\\
(x-2)^2\le0&x=2
\end{array}
\]
\textbf{（2）$\Delta<0$。}二次式恒为正。例如
\[x^2+2x+3=(x+1)^2+2>0.\]
所以它大于（或等于）$0$ 时，解为所有实数；小于（或等于）$0$ 时，无解。

一般地，$a>0,\Delta<0$ 时，
\[ax^2+bx+c=a\left(x+\frac b{2a}\right)^2-\frac\Delta{4a}>0.\]
\textbf{注意：}对应方程无实根，并不意味着不等式无解。
\topic{练习与解答}
（1）$5-3x<11$。

解：$-3x<6$，所以 $x>-2$。

（2）$-x^2+x+6\ge0$。

解：化为 $x^2-x-6\le0$，解集为 $[-2,3]$。

（3）$x^2+1\le0$。

解：因为 $x^2+1\ge1>0$，所以无解。
\end{multicols*}
\end{document}
